\section{Problem statement}

With the perception core component implemented and ready to be tested, the
remaining part of the autonomous system to be implemented is naturally the
control logic. Without it, a drone able to detect the center of the nearest
gate in its field of view is nothing but an inanimate object.

Since this application is for drone racing, the expected flight behavior would
be aggressive, fast but accurate, and capable of taking sharp turns without
hitting any obstacles. However, the control and trajectory planning of UAVs can
be an entire research topic of its own, and therefore cannot be included in the
scope of this project.\\

Initially, an elaborated control algorithm based on reinforcement learning was
planned, but since the focus of this work is on the perception, and by lack of
time, the fallback alternative is a PID-based controller. Naturally, drone
racing cannot be resumed to flying towards a detected obstacle, and even though
PID controllers are good at maintaining a desired position, an algorithm on top
of it is need to supervise the different states and actions present in the
decision making process.

To do so, a state machine is designed and implemented to run on-board, and to
take decisions depending on the state the UAV is in, as well as the desired
outcomes. The simplicity and efficiency of state machines has been proven
before, notably in autonomous drone racing by TU Delft~\etal~\cite{State
Machine based High Level Navigation}. \todo{Find the paper}
Using this approach, coupled with a PID controller to generate velocity
commands to the lower-level controller, it makes it easy and time effective to
test the gate center detection models under different conditions.\\

The controller is implemented in C++ as a ROS node. As such, it can easily be
interfaced with the perception, which is a Python ROS node, and with the safety
system which acts as a proxy between this controller and the lower-level
controller.

\begin{figure}[h]
	\centering
	\includegraphics[width=\textwidth]{figure/300x300.png}
	\caption{Gate center alginment via PID control.}
\end{figure}

As stated previously, the control of the drone is reduced to aligning the image
center of the camera -- technically the drone itself -- with the center of the
predicted region window. The PID controller is therefore calculating the error
as the difference, in pixels , between the image center at $(\frac{W}{2},
\frac{H}{2})$ (for the width and height of the image) and the center of the
predicted window.

To make things simpler and more efficient, the UAV's alignment is achieved by
giving it a constant forward velocity, while the controller sets its upward and
yaw velocities in real time. By choosing the yaw rotation over a horizontal
movement, the drone becomes more agile and faster to approach a gate.
